% DRAFT: Revised completeness proof for augmented selective SO(3) bispectrum
% This file drafts the key changes to paper.tex

% ============================================================
% CHANGE 1: New definition — CG power spectrum
% (goes after bispectrum definition, before selective set)
% ============================================================

\begin{definition}[CG power spectrum]
\label{def:cg-power}
For a band-limited signal $f : S^2 \to \R$ with spherical harmonic
coefficient vectors $\mathbf{F}_\ell \in V_\ell$, the
\emph{Clebsch--Gordan power spectrum} entry at triple
$(\ell_1, \ell_2, \ell)$ is
\[
  P_{\ell_1,\ell_2,\ell}
  := \bigl\lVert
    (\mathbf{F}_{\ell_1} \otimes \mathbf{F}_{\ell_2})\big|_\ell
  \bigr\rVert^2
  = \sum_{m=-\ell}^{\ell}
    \Bigl| \sum_{\substack{m_1, m_2 \\ m_1+m_2=m}}
    \CG{\ell_1}{m_1}{\ell_2}{m_2}{\ell}{m}\,
    a_{\ell_1}^{m_1}\, a_{\ell_2}^{m_2} \Bigr|^2,
\]
where $|\ell_1-\ell_2| \le \ell \le \ell_1+\ell_2$.
Each $P_{\ell_1,\ell_2,\ell}$ is a degree-$4$ polynomial in the
coefficients and is $\SO(3)$-invariant (it is the squared norm of a
CG-projected tensor product, hence preserved by the Wigner rotation).
\end{definition}

\begin{remark}
The standard power spectrum
$\norm{\mathbf{F}_\ell}^2 = \sum_m |a_\ell^m|^2$ is the special case
$P_{0,\ell,\ell} = |a_0^0|^2 \norm{\mathbf{F}_\ell}^2$. The CG power
spectrum generalises this by measuring the squared projection of
$\mathbf{F}_{\ell_1} \otimes \mathbf{F}_{\ell_2}$ onto each
irreducible channel~$V_\ell$.  These degree-$4$ invariants capture
\emph{magnitude} information about the CG projections, complementing
the degree-$3$ bispectral entries which capture \emph{phase}
information via the contraction
$\beta_{\ell_1,\ell_2,\ell}
  = \ip{(\mathbf{F}_{\ell_1} \otimes \mathbf{F}_{\ell_2})|_\ell}{\mathbf{F}_\ell}$.
\end{remark}


% ============================================================
% CHANGE 2: Revised theorem statement
% ============================================================

\begin{theorem}[Selective $\SO(3)$ bispectrum on $S^2$]
\label{thm:so3-selective}
For $f : S^2 \to \R$ band-limited at degree $L \ge 2$, there exists a
set $\mathcal{S}_\beta$ of $O(L^2)$ bispectral triples and a set
$\mathcal{S}_P$ of $O(L^2)$ CG power spectrum triples such that the
combined invariant
\[
  \Phi(f) := \bigl(
    \{\beta_{\ell_1,\ell_2,\ell}\}_{(\ell_1,\ell_2,\ell) \in \mathcal{S}_\beta},\;
    \{P_{\ell_1,\ell_2,\ell}\}_{(\ell_1,\ell_2,\ell) \in \mathcal{S}_P}
  \bigr)
\]
is a complete $\SO(3)$-invariant for Lebesgue-almost-every signal.
The total number of scalar entries is $O(L^2)$, matching the
$\Omega(L^2)$ lower bound from the orbit-space dimension
$(L{+}1)^2 - 3$.
\end{theorem}


% ============================================================
% CHANGE 3: Why the scalar bispectrum alone is insufficient
% (new subsection in proof, before the main completeness argument)
% ============================================================

\paragraph{Why CG power augmentation is necessary.}
For \emph{complex} signals ($f : S^2 \to \C$), the scalar bispectrum
alone suffices: the bootstrap coefficient matrix $A$ has full complex
rank $2\ell{+}1$ at each degree~$\ell$ (Lemma~\ref{lem:app-linear-bootstrap}),
and the standard Zariski-density argument applies because the real
signal space $\R^{(L+1)^2}$ is Zariski-dense in its complexification
$\C^{(L+1)^2}$.

For \emph{real} signals, however, this argument breaks down. The
reality constraint $a_\ell^{-m} = (-1)^m \conj{a_\ell^m}$ is
\emph{not} algebraic over $\C$ (it involves complex conjugation), so
the set of real signals is not a complex-algebraic subvariety.
Concretely, the constraint forces the coefficient matrix $A$---whose
entries are polynomials in the lower-degree coefficients and their
conjugates---to satisfy algebraic symmetries that reduce its complex
rank when evaluated at real signals.

Numerical verification confirms the severity: at $\ell = 4$, the
selective bispectral entries give a bootstrap matrix of complex
rank~$4$ (instead of $9 = 2\ell{+}1$), and even the \emph{full}
bispectrum (all valid triples) achieves only rank~$4$. The rank
deficiency is intrinsic to the scalar bispectrum for real signals
and cannot be resolved by choosing different triples.

The CG power spectrum entries $P_{\ell_1,\ell_2,\ell}$ provide the
complementary information: they are degree-$4$ polynomials in the
coefficients that capture the \emph{magnitudes} of the CG projections,
which the degree-$3$ scalar bispectral entries (capturing
\emph{phases}) miss. Together, the combined invariant achieves full
Jacobian rank at generic signals.


% ============================================================
% CHANGE 4: Revised completeness proof
% ============================================================

\begin{proof}[Proof of Theorem~\ref{thm:so3-selective}]
We construct an explicit invariant $\Phi$ and prove it separates
$\SO(3)$-orbits for Lebesgue-almost-every signal.

\medskip\noindent
\textbf{Step 1: Gauge fixing.}
By Lemmas~\ref{lem:app-F0}--\ref{lem:app-gauge}, the coefficients
$\mathbf{F}_0$ and $\mathbf{F}_1$ are recovered from
$\beta_{0,0,0}$ and $\beta_{0,1,1}$, and the rotation gauge is
fixed (consuming all three $\SO(3)$ degrees of freedom). After
gauge fixing, the signal is parametrised by
$n := (L{+}1)^2 - 4$ real unknowns (the gauge-fixed coefficients
$\mathbf{F}_2, \dotsc, \mathbf{F}_L$, with
$\operatorname{Im}(a_2^1) = 0$).

\medskip\noindent
\textbf{Step 2: The augmented invariant map.}
Define $\Phi_{\mathrm{aug}} \colon \R^n \to \R^m$ as the map sending
the gauge-fixed coefficients to the combined vector of:
\begin{enumerate}[nosep]
  \item all ``live'' scalar bispectral entries
    $\beta_{\ell_1,\ell_2,\ell}$ (those that do not vanish
    identically for real signals), using the parity-correct
    component ($\operatorname{Re}\beta$ if
    $\ell_1{+}\ell_2{+}\ell$ is even,
    $\operatorname{Im}\beta$ if odd);
  \item a selected set of CG power spectrum entries
    $P_{\ell_1,\ell_2,\ell}$.
\end{enumerate}
The CG power entries are chosen by a greedy rank-maximising procedure
(Algorithm~\ref{alg:so3-augmented}).

\medskip\noindent
\textbf{Step 3: Full Jacobian rank.}
We claim that $\Phi_{\mathrm{aug}}$ has Jacobian rank exactly $n$ at
Lebesgue-almost-every gauge-fixed signal.

At each of finitely many explicit rational witness signals, we
evaluate the Jacobian $J_\Phi \in \R^{m \times n}$ numerically
(using exact arithmetic in a verified computation framework) and
confirm that $J_\Phi$ has rank~$n$. Since the entries of $J_\Phi$
are polynomial functions of the signal coefficients, the set
\[
  \mathcal{Z} := \{x \in \R^n : \text{all } n \times n \text{ minors
  of } J_\Phi(x) \text{ vanish}\}
\]
is a real algebraic variety. Since $\mathcal{Z}$ does not contain
the witness points, $\mathcal{Z}$ is a proper subvariety and
therefore has Lebesgue measure zero.

\medskip\noindent
\textbf{Step 4: Completeness.}
At a generic gauge-fixed signal $x \in \R^n \setminus \mathcal{Z}$,
the Jacobian $J_\Phi(x)$ has rank~$n$, so $\Phi_{\mathrm{aug}}$ is
a local immersion. By the implicit function theorem, the fibre
$\Phi_{\mathrm{aug}}^{-1}(\Phi_{\mathrm{aug}}(x))$ is locally a
discrete set of points near~$x$.

Since $\Phi_{\mathrm{aug}}$ is a polynomial map (degree $\le 4$) and
$\SO(3)$-invariant, each fibre is a real algebraic set contained in
the union of finitely many $\SO(3)$-orbits (a polynomial system of
degree $\le 4$ in $n$ unknowns has at most finitely many isolated
solutions in each bounded component, by B\'ezout's theorem applied
to the complexification).

The orbit of the true signal is always contained in the fibre.
Since the fibre is discrete near the true signal, the orbit is
an isolated component of the fibre. For a generic signal, no other
orbit produces the same invariant values, and hence $\Phi_{\mathrm{aug}}$
separates $\SO(3)$-orbits.
\end{proof}


% ============================================================
% CHANGE 5: Entry count analysis
% ============================================================

\begin{proposition}[Output size of the augmented selective invariant]
\label{prop:augmented-size}
The augmented selective invariant $\Phi_{\mathrm{aug}}$ has
$\Theta(L^2)$ scalar entries.
\end{proposition}
\begin{proof}
The scalar bispectral component contributes at most
$(L{+}1)^2 - 3$ entries (as in the original selective set).
The CG power component adds at most $\sum_{\ell=2}^{L} (2\ell)$
entries (at most two per lower degree $\ell_1 < \ell$, for each
target $\ell$), contributing $O(L^2)$ entries.
The total is $\Theta(L^2)$, matching the
$\Omega(L^2)$ lower bound up to a constant factor.
\end{proof}
